% !PRINTED 290
\weberpage{290}
%<<CONT
or
\[
\begin{aligned}
p\alpha' + c'\beta' &= p\alpha, & \beta' &= p\beta\\
p\gamma' + c'\delta' &= c\alpha+\gamma, & \delta' &= c\beta+\delta.
\end{aligned}
\]
From this there follows, by the third equation: $c'\delta' \equiv
c\alpha+\gamma$, and with the help of the fourth:
\[
c' \equiv \frac{c\alpha+\gamma}{c\beta+\delta} \quad (\mathrm{mod}\ p).
\]

If $c\beta+\delta \equiv 0$, then $c' \equiv \infty$, and if
$c \equiv \infty$, then $c' \equiv \alpha/\beta$.

In the application to the determination of $c$, $c'$ one may
replace the $\alpha$, $\beta$, $\gamma$, $\delta$ by congruent numbers $a$, $b$, $c$, $d$
whose determinant $ad-bc$ is a quadratic residue of $c$, % [sic] German source prints c; mathematically p is expected.
and it can
also be shown that in this way every substitution of the group
$\mathfrak{L}_0$ can be obtained. Every function of $v_c$, then, which
remains unchanged under the substitutions of the group $\mathfrak{L}_0$,
remains unchanged also when a linear substitution is applied to
$\omega$, and is consequently a rational function of $j(\omega)$. But of
the nature of the coefficients in this function this consideration teaches
us nothing, and therefore $\mathfrak{L}_0$ is recognised only as the monodromy
group of the invariant equation.

\begin{center}
\S\,79. \gesperrt{Investigation of the Group $\mathfrak{L}_0$}\textsuperscript{1}).
\end{center}

In the 10th chapter of Band II we investigated the congruence group
and its divisors quite generally. We carry this investigation through
here once more in a special form, so far as it bears on the transformation
problem.

The substitutions contained in $\mathfrak{L}_0$
\weq{1}{\begin{pmatrix}a, & b \\ c, & \partial\end{pmatrix} = \left(z, \frac{-c+az}{\partial-bz}\right),}
in which
\weq{2}{\varDelta = a\partial - bc}
is a quadratic residue of $p$, can, as already remarked above,
be brought into the form:
\weq{3}{\begin{pmatrix}\alpha, & \beta \\ \gamma, & \delta\end{pmatrix} = \left(z, \frac{-\gamma+\alpha z}{\delta-\beta z}\right),}
%CONT>>

\begingroup
\renewcommand{\thefootnote}{1)}
\footnotetext[1]{On the group of linear substitutions $\mathfrak{L}_0$ cf.:
\gesperrt{Galois}, Liouvilles Journal, Bd.\,XI. \gesperrt{Serret},
Algèbre supérieure, Section IV, Chapitre IV. C. \gesperrt{Jordan}, Traité des
Substitutions. \gesperrt{Gierster}, Gruppe der Modulargleichungen.
Mathematische Annalen, Bd.\,18.}
\endgroup
